\chapter{TESTING}

Testing is a crucial part of software development to ensure that the Medora Hospital Management System functions correctly, securely, and efficiently. Proper testing helps identify and eliminate bugs, validate system behavior, and ensure smooth interaction between different hospital modules such as Admin, Reception, Doctor, Pharmacy, Laboratory, and Billing.

The system is tested using different testing methods including Unit Testing and Integration Testing.

\section{Unit Testing}

Unit testing focuses on verifying individual components of the system such as user registration, appointment booking, prescription creation, and billing calculation to ensure they work independently without errors.

\renewcommand{\arraystretch}{1.3}

\begin{longtable}{|p{1.5cm}|p{2.5cm}|p{4cm}|p{4cm}|c|}
\hline
\textbf{Test Case ID} & \textbf{Component} & \textbf{Test Description} & \textbf{Expected Result} & \textbf{Status} \\
\hline
\endfirsthead

\hline
\textbf{Test Case ID} & \textbf{Component} & \textbf{Test Description} & \textbf{Expected Result} & \textbf{Status} \\
\hline
\endhead

\hline
\caption{Unit Test Cases for Medora Hospital Management System} \label{tab:unit_tests} \\
\endlastfoot

UT01 & User Registration & Verify patient registration form validation & User account created successfully & Pass \\
\hline
UT02 & Login Authentication & Check valid and invalid login attempts & Access granted only with correct credentials & Pass \\
\hline
UT03 & Appointment Booking & Validate appointment slot availability & Appointment booked only if slot available & Pass \\
\hline
UT04 & Prescription Creation & Verify doctor can create prescription & Prescription saved successfully & Pass \\
\hline
UT05 & Lab Test Request & Check lab test request generation & Lab test linked to patient correctly & Pass \\
\hline
UT06 & Billing Calculation & Validate total amount calculation & Correct total amount generated & Pass \\
\hline
UT07 & Medicine Stock Update & Verify stock deduction after dispensing & Stock reduced correctly & Pass \\
\hline
UT08 & Role Authorization & Check role-based access restrictions & Unauthorized access prevented & Pass \\
\end{longtable}

\newpage
\section{Integration Testing}

Integration testing ensures that different modules interact properly with each other without errors.

\begin{table}[H]
\centering
\renewcommand{\arraystretch}{1.3}
\small
\begin{tabular}{|c|p{3cm}|p{3cm}|p{3cm}|c|}
\hline
\textbf{Test Case ID} & \textbf{Module Interaction} & \textbf{Test Description} & \textbf{Expected Result} & \textbf{Status} \\
\hline
IT01 & Patient–Reception & Verify patient registration and appointment booking flow & Patient appears in appointment list & Pass \\
\hline
IT02 & Doctor–Pharmacy & Ensure prescription links to patient correctly & Prescription visible in pharmacy & Pass \\
\hline
IT03 & Pharmacy–Billing & Verify billing generated after medicine dispensing & Bill created successfully & Pass \\
\hline
IT04 & Doctor–Lab & Ensure lab test requests appear for lab staff & Lab staff can view test request & Pass \\
\hline
IT05 & Lab–Patient & Verify lab result upload visibility to patient & Patient can view lab report & Pass \\
\hline
IT06 & Patient–Pharmacy (Direct Purchase) & Check patient can upload external prescription & Pharmacist can verify uploaded prescription & Pass \\
\hline
\end{tabular}
\caption{Integration Test Cases for Medora Hospital Management System}
\label{tab:integration_tests}
\end{table}